% Source: https://artofproblemsolving.com/community/c3h1533960_divisibilty_by_3_integers
% Stable identity: aops-topic-1533960-q01
The three conditions all say the same thing about $2n$:
\begin{itemize}
\item $15\mid n$ gives $2n\equiv0\equiv15\pmod{15}$;
\item $17\mid n+1$ gives $2n\equiv-2\equiv15\pmod{17}$;
\item $19\mid n+2$ gives $2n\equiv-4\equiv15\pmod{19}$.
\end{itemize}
Since $15$, $17$, $19$ are pairwise coprime (and $2$ is invertible modulo each), the three conditions together are equivalent to
\[
2n\equiv15\pmod{15\cdot17\cdot19},\qquad\text{i.e.}\qquad 2n-15=4845k
\]
for some integer $k$. For $n$ to be an integer, $k$ must be odd, and $n>0$ needs $k\ge1$. The smallest choice $k=1$ gives $2n=4860$, so $n=\boxed{2430}$.

Check: $2430=15\cdot162$, $2431=17\cdot143$, $2432=19\cdot128$.
